\section{Limitations}
\label{sec:limitations}

The model answers a narrow question --- what an imposed automation and productivity path does to UK output, factor shares and the fiscal aggregates over \genStartYear--\genEndYear{} --- and the boundaries are sharp. They are set out here rather than distributed through the text.

\paragraph{The labour-share result is close to assumed.} At $\varepsilon = \genEpsilon{}$ the labour share is exactly $1-\gamma_t$ in every period. The \genLabourShareFallPp{} percentage-point fall reported in Section~\ref{sec:results-wages} is therefore the calibration target restated, not a model output, and the smooth path from \genBaseLabourSharePct{} to \genShareShiftTfpLabourShareThirty{} per cent demonstrates only that the time-varying exponent works. A model in which the labour share is an equilibrium object requires an elasticity below one, a task structure, or both.

\paragraph{The elasticity of substitution is the single most consequential assumption.} This paper uses $\varepsilon = \genEpsilon{}$, OG-UK's single-sector default. \citet{korinek2026} assume 0.5, and \citet{mollimas2026} argue for approximately 0.2. At a unit elasticity capital deepening cannot move factor shares at all; at an elasticity below one it can, and the mapping from automation intensity to output and wages changes. No sensitivity over $\varepsilon$ is run here, because OG-Core indexes $\varepsilon$ by industry only --- the same limitation the patch of Section~\ref{sec:patch} removes for $\gamma$, at roughly \genPatchEpsilonSites{} sites in the same file.

\paragraph{No unemployment or search frictions.} The model has no labour-market block, so displacement cannot be represented. The OBR's rise in equilibrium unemployment to 5.5~per cent and the \genKorinekTableUnemploymentPct{} per cent of \citet{korinek2026} have no counterpart, and the paper says nothing about who loses work. Aggregate labour input falls (\genShareShiftTfpLabourGapThirtyAbs{} and \genShareShiftLabourGapThirtyAbs{} per cent by \genEndYear{}) through cohort labour-supply choices at a lower wage, which is a different mechanism from involuntary job loss and should not be read as one.

\paragraph{No task reinstatement.} \citet{acemoglu2018} and \citet{acemoglu2022tasks} identify the creation of new labour-intensive tasks as the main countervailing force to automation, and \citet{korinek2026} carry an explicit reinstatement parameter. There is none here, so the labour share cannot recover and the scenarios are, in that specific sense, one-directional by construction.

\paragraph{No adoption or diffusion channel.} The productivity gain arrives on the assumed schedule because nothing in the model can delay it. \citet{korinek2026} carry an adoption parameter that builds over the window, and \citet{karger2026} examine how much of such a gain is measurable so far. % TODO verify characterisation of karger2026
A model without diffusion cannot produce a slower-than-assumed path.

\paragraph{No research and development, and no ideas channel.} Technology is exogenous. The mechanism by which AI raises the productivity of research itself, the subject of \citet{jonestonetti2026}, is absent, so the model cannot generate accelerating growth and cannot speak to the upper end of the published range.

\paragraph{One sector.} All runs are single-sector, so the compositional channels --- different automation intensities and elasticities across industries, and shifts in the composition of demand --- are absent. A multi-sector build exists in the wider OG-UK work but is not used here, and its transition path was not reliable over longer horizons.

\paragraph{The automation-only arm is not a reproduction of the OBR's scenario, and its labour-share target is borrowed.} The OBR publishes no labour-share number for its technological-displacement scenario. The automation-only arm therefore imposes the same \genLabourShareFallPp{}-point intensity as the productivity arm and lets $Z$ deliver the output-neutrality condition that defines the OBR case. Only the neutrality condition is motivated by the OBR; the automation intensity comes from \citet{korinek2026}. The arm's internal name was changed from \texttt{obr} to \texttt{zero\_tfp} to stop the label implying a reproduction of Box 2.2, whose quantitative result (GDP level unchanged) this arm does not deliver.

\paragraph{Output neutrality holds only at fixed inputs.} The automation-only arm's GDP ends \genShareShiftGdpGapThirtyAbs{} per cent above baseline despite being calibrated to leave output unchanged, because the neutrality condition binds at the baseline's fixed $K_0$ and $L_0$ while saving is endogenous. Any comparison with the OBR's ``level of GDP unchanged'' must carry this qualification.

\paragraph{The solved $Z$ levels have no economic meaning.} Both solved values ($\genZShareShiftTfp{}$ and $\genZShareShift{}$) lie below~1 because they offset the \genGammaOnlyGainPct{} per cent output change that moving $\gamma$ alone produces at the baseline's factor levels, which in turn reflects the units of $K$ and $L$ in the model. Only the joint $(\gamma, Z)$ move is interpretable.

\paragraph{Solution checks were disabled, and \genStartYear{} carries numerical error.} OG-Core's \texttt{ENFORCE\_SOLUTION\_CHECKS} was set to \genEnforceChecks{}, and the diagnostics are reported instead (Appendix~\ref{app:diagnostics}). Interior resource-constraint errors are small (\genBaseRcMaxInterior{} at the maximum) and Euler errors are negligible (\genBaseEulerMax{}), but the \genStartYear{} resource-constraint error is \genBaseRcTwentySixPctGdp{} per cent of \genStartYear{} GDP. Levels and ratios in \genStartYear{} carry it. Reported gaps against the baseline largely net it out, because every arm carries an error of the same order, but \genStartYear{} figures should be treated as the least reliable in every table.

\paragraph{Investment is a residual.} Investment is the residual in the resource constraint and absorbs that numerical error, so its gaps are the least reliable series reported. The capital stock is the series on which the output decomposition rests.

\paragraph{The fiscal-rule switch was moved.} The switch of government spending to debt targeting is set to $t_{G1} = \genTGone{}$ (\genTGoneYear{}), past the reported window, so that no fiscal rule binds inside it. This is a deliberate choice to isolate the AI channel; it also means the runs say nothing about how a binding fiscal rule would interact with an AI shock.

\paragraph{No step-arm comparison in the committed results.} The runner can produce stepped and ramped arms in one pass, but the committed results contain the ramped arms only. The step-versus-ramp contrast in Sections~\ref{sec:patch} and~\ref{sec:discussion} is therefore argued from the structure of the model and from the patch's verification run, not from a paired step and ramp in the same results file.

\paragraph{The patch is not upstream.} \texttt{patches/ogcore-0.17.0-firm-gamma-tv.diff} is applied to the installed OG-Core \genOgcoreVersion{} in the run environment. OG-Core itself is unmodified and no pull request has been merged, so these results are not reproducible from a stock installation without applying the diff.

\paragraph{One calibration, not a distribution.} Each scenario is a single deterministic run. There is no uncertainty quantification over $\gamma$, $Z$, $\varepsilon$, the demographic projection or the tax functions, so no interval accompanies any number in this paper.
